\chapter{Introdução}\label{cap:introducao}

\section{Finalidade e problema}
Uma especificação de requisitos de software (SRS) estabelece comportamentos,
interfaces, regras e condições de qualidade observáveis de um produto. Ela
permite que usuários, desenvolvimento e testes trabalhem com o mesmo contrato.
Este guia-modelo conserva a base didática de \citeonline{srstex}, usa
abn\TeX2\ \cite{abntex2} e demonstra uma especificação para um sistema fictício
de inscrições em eventos.

O problema do exemplo é controlar vagas e inscrições sem duplicidade, permitir
que a pessoa acompanhe seus registros e respeitar o prazo de cada evento.
As fichas especificam esse produto; as explicações orientam a adaptação do modelo.
Os objetivos e metas apresentados são escolhas do exemplo, sem constituir
resultados de entrevistas ou testes de um software implementado.

\section{Objetivos do produto}
\begin{quadrosrs}
\caption{Objetivos e comportamentos associados}\label{tab:objetivos}
\begin{tabelasrs}{L{0.14\textwidth}YY}
\cabecalhosrs\textbf{ID} & \textbf{Objetivo} & \textbf{Comportamento observável}\\\midrule
OBJ-001 & Encontrar eventos disponíveis. & Catálogo filtra publicação e início e apresenta vagas.\\
OBJ-002 & Realizar e acompanhar inscrições válidas. & Confirmação persistida, controle de capacidade e consulta pessoal.\\
OBJ-003 & Desistir de uma inscrição no prazo. & Cancelamento pelo titular e liberação de uma única vaga.\\
\end{tabelasrs}
\fonte{Elaboração didática deste guia.}
\end{quadrosrs}

\Needspace{17\baselineskip}
\section{Escopo e fronteiras}
O produto consome eventos de um catálogo externo e identidades de um provedor
externo. É responsável pelo catálogo público de consulta e pelo ciclo de vida
das inscrições. A fronteira separa essas funções da administração dos serviços
que fornecem eventos e autenticação.

\begin{quadrosrs}
\caption{Delimitação do escopo}
\begin{tabelasrs}{YY}
\cabecalhosrs\textbf{Dentro do produto} & \textbf{Fora do produto}\\\midrule
Consultar eventos e inscrições pessoais. & Publicar e editar eventos no serviço de origem.\\
Criar e cancelar inscrições; controlar vagas. & Cobrar pagamentos e emitir certificados.\\
Aplicar sessão e titularidade nas operações. & Cadastrar credenciais e administrar o provedor de identidade.\\
\end{tabelasrs}
\fonte{Elaboração didática deste guia.}
\end{quadrosrs}

\section{Convenções e priorização}
Os identificadores são estáveis: \texttt{RF} para requisitos funcionais,
\texttt{RNF} para qualidade, \texttt{RN} para regras, \texttt{UC} para casos de uso
 e \texttt{CT} para testes. O número não muda quando uma seção é reorganizada.
``Deve'' estabelece uma obrigação; origem, condições e aceite delimitam sua
interpretação. Cada prioridade é explícita, sem herança automática entre fichas.

\begin{quadrosrs}
\caption{Escala de prioridade}
\begin{tabelasrs}{L{0.17\textwidth}Y}
\cabecalhosrs\textbf{Prioridade} & \textbf{Significado}\\\midrule
Alta & Essencial para uma operação correta ou para proteção dos registros.\\
Média & Qualidade complementar relevante para executar os fluxos previstos.\\
Baixa & Melhoria opcional que não impede os objetivos essenciais.\\
\end{tabelasrs}
\fonte{Elaboração didática deste guia.}
\end{quadrosrs}

Evite ``rápido'', ``seguro'' ou ``intuitivo'' sem condições e medidas. O aceite
segue \textbf{Dado / Quando / Então}: situação inicial, ação e resultado
observável. A engenharia de requisitos também envolve processos ao longo do
ciclo de vida \cite{iso29148}; este guia não certifica conformidade com a norma.

\section{Público e organização}
Responsáveis pelo produto consultam objetivos, contexto e regras. A equipe de
desenvolvimento usa também dados, interfaces e fichas. A equipe de testes usa
os casos de uso e a matriz da \autoref{cap:verificacao}. O glossário e o modelo
de estados estão no \autoref{ap:glossario}. A ordem de leitura acompanha o
caminho necessidade, regra, contrato, requisito e verificação.
